% Inhalt beider Fassungen (graufieber.tex: Blattfassung, graufieber-buch.tex: Buchfassung).
% Die Dateien in build/ erzeugt das Build-Skript.
\fwtitle{Graufieber}{Ein Abenteuer für The One Ring, 2nd Edition\fwsubsep Tharbad und Swanfleet, Spätsommer TA 2970}

\input{build/body.tex}

% ---- Rechtehinweis am Ende des Textes
\fwnotice

% ---- Karten für den Loremaster (nummeriert, quer)
\fwlmmap[Karten für den Loremaster]{lm-weidenhof}{Karte: Der Weidenhof}%
  {Loremaster-Version · Nummern wie in Teil 3}{assets/maps/map-weidenhof-de.jpg}{Der Weidenhof}
\fwlmmap{lm-reiher}{Karte: Die Reiherinsel}%
  {Loremaster-Version · Nummern wie im Höhepunkt}{assets/maps/map-reiherinsel-de.jpg}{Die Reiherinsel}
\label{fw-lastnum}

% ---- Handouts für die Spieler (ohne Seitenzahl)
\input{build/letters.tex}
\fwhandoutletter{h-calenmir}{Calenmirs Brief}{calenmir}{C.}
\fwhandoutmap[Handouts für die Spieler]{h-tharbad}{assets/maps/map-tharbad-players-de.jpg}{Karte: Tharbad und die Swanfleet}
\fwhandoutmap{h-weidenhof}{assets/maps/map-weidenhof-players-de.jpg}{Karte des Weidenhofs}
\fwhandoutletter{h-vorlund}{Vorlunds Bericht}{vorlund}{Euer Diener V.}
\fwhandoutmap{h-reiher}{assets/maps/map-reiherinsel-players-de.jpg}{Karte der Reiherinsel}
